\documentclass[a4paper,12pt]{article}
\usepackage[utf8]{inputenc}
\usepackage[T1]{fontenc}
\usepackage[french]{babel}
\usepackage{geometry}
\usepackage{enumitem}
\usepackage{xcolor}
\usepackage{hyperref}

\geometry{margin=2.5cm}

\title{Rapport de Revue de Code : Analyse des Fichiers Orphelins et Obsolètes}
\author{Cline AI}
\date{\today}

\begin{document}

\maketitle

\section{Introduction}
L'objectif de ce rapport est d'identifier les fichiers et dossiers orphelins, obsolètes ou redondants au sein du projet. L'analyse a été effectuée en comparant l'arborescence complète du système de fichiers avec les configurations actives, notamment le fichier \texttt{docker-compose.yml} et les structures d'applications.

\section{Analyse des Dossiers d'Archives et de Sauvegardes}
Plusieurs dossiers ont été identifiés comme ne faisant pas partie du flux de production et servant uniquement de stockage temporaire ou de sauvegarde.

\begin{itemize}
    \item \textbf{\texttt{archive\_refactor/}} : Contient des versions antérieures de scripts, des dumps SQL et des configurations de refactorisation.
    \item \textbf{\texttt{backups/}} : Contient des archives compressées (\texttt{.tar.gz}) de sauvegardes locales.
    \item \textbf{\texttt{trash/}} : Dossier de rebuts contenant explicitement des fichiers supprimés.
    \item \textbf{\texttt{remote\_ai\_backup/}} : Sauvegardes de scripts liés à l'IA.
\end{itemize}

\section{Fichiers de Configuration Redondants}
Le projet présente une multiplication des fichiers de configuration, ce qui peut mener à des erreurs de déploiement.

\subsection{Docker Compose}
Seul \texttt{docker-compose.yml} est utilisé comme configuration principale. Les fichiers suivants sont considérés comme obsolètes ou des variantes de test :
\begin{itemize}
    \item \texttt{docker-compose.ai.yml}
    \item \texttt{docker-compose.dev.yml}
    \item \texttt{docker-compose.planka.yml}
    \item \texttt{docker-compose.scodoc-test.yml}
    \item \texttt{docker-compose.yml.final\_https}
    \item \texttt{docker-compose.yml.native\_https}
    \item \texttt{docker-compose.yml.parallel\_test}
    \item \texttt{docker-compose.yml.restore}
\end{itemize}

\subsection{Configuration Nginx et Environnement}
\begin{itemize}
    \item \texttt{nginx.conf.native\_https} et \texttt{nginx.conf.restore} : Redondants face à la configuration spécifique du dashboard.
    \item \texttt{.env.ai} et \texttt{.env.backup\_demo\_20260205} : Versions obsolètes des variables d'environnement.
\end{itemize}

\section{Scripts de Test et Debug}
Plusieurs scripts Python de test se trouvent à la racine sans être intégrés dans un cycle de CI/CD ou référencés dans les services :
\begin{itemize}
    \item \texttt{debug\_scodoc\_login.py}
    \item \texttt{scodoc\_tester.py}
    \item \texttt{test\_scodoc\_web\_v2.py}
    \item \texttt{test\_scodoc\_web.py}
    \item \texttt{test\_scodoc.py}
\end{itemize}

\section{Artefacts de Compilation}
Le dossier \texttt{tuto-SITC/} et les documents de sauvegarde contiennent des fichiers générés par la compilation \LaTeX{} qui ne devraient pas être suivis par le contrôle de version :
\begin{itemize}
    \item Fichiers \texttt{.aux}, \texttt{.fdb\_latexmk}, \texttt{.fls}, \texttt{.out}, \texttt{.synctex.gz}, \texttt{.toc}.
\end{itemize}

\section{Autres Fichiers Orphelins}
\begin{itemize}
    \item \texttt{server.php} : Fichier isolé à la racine sans lien avec les services actuels.
    \item \texttt{go.sh} et \texttt{gogo.sh} : Scripts de lancement redondants face à \texttt{start.sh}.
\end{itemize}

\section{Conclusion et Recommandations}
Il est fortement recommandé de :
\begin{enumerate}
    \item Déplacer les dossiers \texttt{archive\_refactor}, \texttt{backups} et \texttt{trash} en dehors du répertoire de travail ou les compresser dans une archive externe.
    \item Nettoyer les variantes de \texttt{docker-compose.yml} pour ne garder que la version active.
    \item Créer un fichier \texttt{.gitignore} plus exhaustif pour exclure les artefacts \LaTeX{} et les logs (\texttt{.log}).
    \item Centraliser les scripts de test dans un dossier \texttt{tests/} dédié.
\end{enumerate}

\end{document}
</write_to_file>
<task_progress>
- [x] Explorer l'arborescence complète du projet
- [x] Analyser les fichiers pour détecter les orphelins/obsolètes
- [x] Identifier les dossiers d'archives et de backups
- [x] Vérifier les références croisées (imports, config)
- [x] Rédiger le rapport final en LaTeX
</task_progress>